\documentclass[12pt]{article}

\usepackage{amsmath}
\usepackage{amssymb}
\usepackage{physics}
\usepackage{enumitem}
\usepackage{tikz}
\usetikzlibrary{arrows.meta,calc,angles,quotes,decorations.pathmorphing}
\usepackage{geometry}
\geometry{margin=1in}

\setlist[enumerate,1]{label=\textbf{Q\arabic*.}}
\setlist[enumerate,2]{label=(\Alph*)}

\begin{document}

\begin{center}
\Large\textbf{SET 15 (Upgraded – High Difficulty, Syllabus Strict)}
\end{center}

\begin{enumerate}

%====================================================
\item A small bead of mass $m$ slides without friction on a circular hoop of radius $R$ fixed in a vertical plane. The hoop rotates with constant angular speed $\omega$ about its vertical diameter. For $\omega^2 > \dfrac{g}{R}$, two symmetric equilibrium positions appear at angle $\theta_0$ from the lowest point.

If $\theta = \theta_0 + \phi$ with $|\phi|\ll1$, determine the angular frequency of small oscillations.

\begin{center}
\begin{tikzpicture}[scale=1]
\draw (0,0) circle (2);
\draw[dashed] (0,-2)--(0,2);
\fill (1.3,1.4) circle (2pt);
\node[right] at (1.3,1.4){$m$};
\draw[->] (0,-2.6)--(0,-3.4);
\node at (0,-3.8){$\omega$};
\end{tikzpicture}
\end{center}

\begin{enumerate}
\item $\sqrt{\omega^2 - \dfrac{g}{R}}$
\item $\sqrt{2\omega^2 - \dfrac{g}{R}}$
\item $\sqrt{2\left(\omega^2 - \dfrac{g}{R}\right)}$
\item $\sqrt{\dfrac{g}{R}}$
\end{enumerate}

%====================================================
\item A uniform solid sphere rolls without slipping down an incline of angle $30^\circ$ from rest. After descending a vertical height of $1.0\,\text{m}$, its speed is

\begin{enumerate}
\item $3.7\,\text{m/s}$
\item $4.2\,\text{m/s}$
\item $3.1\,\text{m/s}$
\item $2.8\,\text{m/s}$
\end{enumerate}

%====================================================
\item A mass $m$ attached to a spring of constant $k$ rotates in a horizontal plane about the fixed end of the spring with angular speed $\Omega$. The equilibrium length becomes $r_0$. For small radial oscillations about $r_0$, the oscillation frequency is

\begin{enumerate}
\item $\sqrt{\dfrac{k}{m} - \Omega^2}$
\item $\sqrt{\dfrac{k}{m} - 3\Omega^2}$
\item $\sqrt{\dfrac{k}{m} + \Omega^2}$
\item $\sqrt{\dfrac{k}{m}}$
\end{enumerate}

%====================================================
\item A solid sphere of radius $R$ carries volume charge density $\rho(r)=\rho_0\left(1-\dfrac{r}{R}\right)$. Using Gauss’s law and proper integration, the electric field at $r=\dfrac{R}{2}$ equals

\begin{enumerate}
\item $\dfrac{\rho_0 R}{24\epsilon_0}$
\item $\dfrac{\rho_0 R}{12\epsilon_0}$
\item $\dfrac{\rho_0 R}{8\epsilon_0}$
\item $\dfrac{\rho_0 R}{6\epsilon_0}$
\end{enumerate}

%====================================================
%====================================================
\item A conducting rod of mass $m$ and resistance $R$ slides without friction on parallel rails separated by distance $l$ in a uniform magnetic field $B$ perpendicular to the plane. The rails are connected by a resistor $R$ at one end.

If the rod is given an initial speed $v_0$, determine the total distance travelled before coming to rest.

\begin{enumerate}
\item $\dfrac{mRv_0}{B^2l^2}$
\item $\dfrac{mRv_0}{2B^2l^2}$
\item $\dfrac{2mRv_0}{B^2l^2}$
\item $\dfrac{mRv_0^2}{B^2l^2}$
\end{enumerate}

%====================================================
\item A parallel plate capacitor connected to a battery is completely filled with dielectric of constant $K$. The battery remains connected. The ratio of final energy stored to initial energy stored is

\begin{enumerate}
\item $K$
\item $K^2$
\item $\dfrac{1}{K}$
\item $\sqrt{K}$
\end{enumerate}

%====================================================
\item A long solenoid of radius $R$ carries time varying current such that magnetic field inside is $B=kt$. The induced electric field at radial distance $r$ from axis is

\begin{enumerate}
\item $\dfrac{kr}{2}$ for $r<R$, and $\dfrac{kR^2}{2r}$ for $r>R$
\item $\dfrac{kR}{2}$ for all $r$
\item $\dfrac{kr}{2}$ for all $r$
\item Zero for $r>R$
\end{enumerate}

%====================================================
\item In Young’s double slit experiment, slit intensities are in ratio $9:4$. For central point to become completely dark, additional phase difference introduced must be

\begin{center}
\begin{tikzpicture}[scale=0.9]
\draw (0,0)--(0,3);
\draw (-0.1,1)--(0.1,1);
\draw (-0.1,2)--(0.1,2);
\draw (4,0)--(4,3);
\draw[dashed] (0,1)--(4,1.5);
\draw[dashed] (0,2)--(4,1.5);
\node at (4.6,1.5){Screen};
\end{tikzpicture}
\end{center}

\begin{enumerate}
\item $\pi$
\item $\dfrac{\pi}{2}$
\item $2\pi$
\item Impossible
\end{enumerate}

%====================================================
\item A thin convex lens of focal length $f$ moves with uniform speed $u \ll c$ along its principal axis away from a stationary object placed at distance $2f$ from the lens.

To first order in $u/c$, the image formed in laboratory frame is

\begin{enumerate}
\item at distance $2f$ and moving away
\item at distance $2f$ and stationary
\item at distance $2f\left(1+\dfrac{u}{c}\right)$
\item at distance $2f\left(1-\dfrac{u}{c}\right)$
\end{enumerate}

%====================================================
\item A thin film of refractive index $1.5$ is illuminated normally by light of wavelength $600\,\text{nm}$ in air. Minimum thickness for first bright reflection is

\begin{enumerate}
\item $100\,\text{nm}$
\item $200\,\text{nm}$
\item $300\,\text{nm}$
\item $400\,\text{nm}$
\end{enumerate}

%====================================================
\item An electromagnetic wave of electric field amplitude $E_0$ is incident normally on a perfectly absorbing surface. The radiation pressure exerted is

\begin{enumerate}
\item $\dfrac{\epsilon_0 E_0^2}{2}$
\item $\epsilon_0 E_0^2$
\item $\dfrac{\epsilon_0 E_0^2}{4}$
\item $\dfrac{\epsilon_0 E_0^2}{c}$
\end{enumerate}

%====================================================
\item Light of wavelength $500\,\text{nm}$ falls on metal of work function $2.0\,\text{eV}$. Using $hc=1240\,\text{eV nm}$, stopping potential is approximately

\begin{enumerate}
\item $0.48\,\text{V}$
\item $0.60\,\text{V}$
\item $0.96\,\text{V}$
\item $1.20\,\text{V}$
\end{enumerate}

%====================================================
\item In Compton scattering, photon of wavelength equal to Compton wavelength $\lambda_c$ scatters at $90^\circ$. The scattered wavelength becomes

\begin{enumerate}
\item $2\lambda_c$
\item $1.5\lambda_c$
\item $\sqrt{2}\lambda_c$
\item $3\lambda_c$
\end{enumerate}

%====================================================
\item A Zener diode is connected in series with resistor to a supply. In breakdown region, small change in supply voltage occurs. Assuming ideal Zener, fractional change in load voltage is

\begin{enumerate}
\item $0$
\item $\dfrac{\Delta V}{V_Z}$
\item $\dfrac{\Delta V}{V_0}$
\item $\dfrac{R_L}{R}\dfrac{\Delta V}{V_0}$
\end{enumerate}

%====================================================
\item A simple pendulum of length $l$ and bob of mass $m$ carries charge $q$. It is placed in uniform horizontal electric field $E$. The equilibrium shifts by angle $\theta_0$. For small oscillations about new equilibrium, the time period is

\begin{center}
\begin{tikzpicture}[scale=1]
\draw (0,0)--(0,-2);
\fill (0,-2) circle (3pt);
\draw[->] (1,-1)--(2,-1);
\node at (2.3,-1){$E$};
\draw[dashed] (0,0)--(0.5,-2);
\end{tikzpicture}
\end{center}

\begin{enumerate}
\item $2\pi\sqrt{\dfrac{l}{g}}$
\item $2\pi\sqrt{\dfrac{l}{\sqrt{g^2+(qE/m)^2}}}$
\item $2\pi\sqrt{\dfrac{l}{g+qE/m}}$
\item $2\pi\sqrt{\dfrac{l}{g-qE/m}}$
\end{enumerate}

\end{enumerate}

\end{document}